\chapter{Defects, Boundaries, and Extended Symmetry Data}
\label{ch:volume-iii-defects-boundaries-extended-symmetry-data}

\section{Conformal Defects}

Let the ambient CFT live on \(\mathbb R^D\).  A flat \(p\)-dimensional
conformal defect is supported on a plane
\[
  \Sigma=\mathbb R^p\subset\mathbb R^D.
\]
It preserves the subgroup
\[
  SO(p+1,1)\times SO(D-p)
  \subset SO(D+1,1),
\]
where \(SO(p+1,1)\) is the conformal group along \(\Sigma\) and
\(SO(D-p)\) rotates the transverse directions.  A defect is denoted
\[
  \mathcal D[\Sigma].
\]
Correlation functions in the presence of the defect are
\[
  \left\langle
    \mathcal O_1(x_1)\cdots\mathcal O_n(x_n)
  \right\rangle_{\mathcal D}
  =
  \frac{
    \left\langle
      \mathcal D[\Sigma]\,
      \mathcal O_1(x_1)\cdots\mathcal O_n(x_n)
    \right\rangle
  }{\langle\mathcal D[\Sigma]\rangle}.
\]

Local operators can live on the defect.  A defect local operator
\(\widehat{\mathcal O}(\sigma)\), with \(\sigma\in\Sigma\), transforms as a
primary under \(SO(p+1,1)\) and in a representation of the transverse rotation
group \(SO(D-p)\).  Its quantum numbers are
\[
  \widehat\Delta,\qquad \widehat\rho,\qquad r_\perp,
\]
where \(\widehat\Delta\) is the defect scaling dimension,
\(\widehat\rho\) is the spin along the defect, and \(r_\perp\) is the
transverse-spin representation.

\begin{figure}[H]
  \centering
  \resizebox{0.92\textwidth}{!}{%
  \begin{tikzpicture}[x=1cm,y=1cm,every node/.style={font=\scriptsize}]
    \tikzset{
      arrow/.style={-{Latex[length=2mm]},line width=0.85pt},
      defect/.style={violet!80!black,line width=1.2pt},
      bulk/.style={circle,fill=blue!65!black,inner sep=2pt},
      dop/.style={circle,fill=orange!85!black,inner sep=1.8pt}
    }
    \draw[defect] (-4.4,0)--(4.4,0);
    \node[violet!80!black] at (3.65,0.32) {defect \(\Sigma\)};
    \node[bulk] at (-1.2,1.45) {};
    \node[blue!65!black] at (-1.2,1.80) {\(\mathcal O(x_\perp,\sigma)\)};
    \draw[arrow,blue!65!black] (-1.2,1.28)--(-1.2,0.20);
    \foreach \x/\lab in {-2.0/{\widehat{\mathcal O}_1},
      -0.9/{\widehat{\mathcal O}_2},0.35/{\widehat{\mathcal O}_3},
      1.65/{\cdots}} {
      \node[dop] at (\x,0) {};
      \node[orange!85!black,below] at (\x,-0.02) {\(\lab\)};
    }
    \node[green!45!black,align=center] at (2.3,1.25)
      {bulk-to-defect OPE\\expands a bulk operator\\into defect primaries};
  \end{tikzpicture}
  }
  \caption{A conformal defect supports its own local operator spectrum, and
  bulk operators admit an expansion into defect primaries.}
  \label{fig:bulk-to-defect-ope}
\end{figure}

\section{Bulk-to-Defect OPE}

Write a point in the ambient space as
\[
  x=(\sigma,x_\perp),
  \qquad
  \sigma\in\mathbb R^p,
  \quad
  x_\perp\in\mathbb R^{D-p}.
\]
For a scalar bulk primary \(\mathcal O\) of dimension \(\Delta\), the
bulk-to-defect OPE has the form
\begin{equation}
  \mathcal O(\sigma,x_\perp)
  =
  \sum_{\widehat{\mathcal O}}
  b_{\mathcal O\widehat{\mathcal O}}\,
  |x_\perp|^{\widehat\Delta-\Delta}
  \mathcal C_{\widehat{\mathcal O}}(x_\perp,\partial_\sigma)
  \widehat{\mathcal O}(\sigma).
  \label{eq:bulk-to-defect-ope}
\end{equation}
The sum runs over defect primaries compatible with the transverse-spin
selection rule.  The differential operator
\(\mathcal C_{\widehat{\mathcal O}}\) generates defect descendants and
transverse tensor structures.  The coefficients
\[
  b_{\mathcal O\widehat{\mathcal O}}
\]
are part of the defect CFT data.

The one-point function of a scalar bulk primary in the presence of a flat
defect is fixed up to a coefficient:
\[
  \left\langle\mathcal O(\sigma,x_\perp)\right\rangle_{\mathcal D}
  =
  \frac{a_{\mathcal O}}{|x_\perp|^\Delta}.
\]
The coefficient \(a_{\mathcal O}\) is the bulk-to-defect OPE coefficient of
\(\mathcal O\) into the defect identity.

\section{Defect Crossing}

Consider a two-point function of identical scalar bulk primaries in the
presence of a defect:
\[
  \left\langle
    \mathcal O(x_1)\mathcal O(x_2)
  \right\rangle_{\mathcal D}.
\]
There are two natural expansions.  In the bulk channel, apply the ordinary
bulk OPE \(\mathcal O\times\mathcal O\), and then take one-point functions in
the defect background.  In the defect channel, apply the bulk-to-defect OPE to
each bulk operator and then use defect two-point functions.  Equality of the
two expansions is defect crossing.

For identical scalar external operators, with one cross ratio \(\xi\), the
crossing equation takes the form
\begin{equation}
  \sum_{\mathcal X}
  \lambda_{\mathcal O\mathcal O\mathcal X}
  a_{\mathcal X}
  G_{\mathcal X}^{\rm bulk}(\xi)
  =
  \sum_{\widehat{\mathcal X}}
  b_{\mathcal O\widehat{\mathcal X}}^2
  G_{\widehat{\mathcal X}}^{\rm defect}(\xi),
  \label{eq:defect-crossing-equation}
\end{equation}
where \(\xi\) denotes the conformal cross ratio appropriate to two bulk
points and a flat defect.  The functions \(G^{\rm bulk}\) and
\(G^{\rm defect}\) are fixed by the preserved conformal group.  Equation
\eqref{eq:defect-crossing-equation} is the defect bootstrap equation.

\begin{figure}[H]
  \centering
  \resizebox{0.92\textwidth}{!}{%
  \begin{tikzpicture}[x=1cm,y=1cm,every node/.style={font=\scriptsize}]
    \tikzset{
      arrow/.style={-{Latex[length=2mm]},line width=0.85pt},
      defect/.style={violet!80!black,line width=1.1pt},
      line/.style={blue!65!black,line width=0.9pt},
      op/.style={circle,fill=blue!65!black,inner sep=2pt}
    }
    \begin{scope}[shift={(-3.3,0)}]
      \draw[defect] (-1.7,-1.0)--(1.7,-1.0);
      \node[op] (a) at (-0.7,0.8) {};
      \node[op] (b) at (0.8,0.7) {};
      \draw[line] (a)--(b) node[midway,above] {bulk OPE};
      \draw[line,dashed] (0.05,0.72)--(0.05,-1.0);
      \node[violet!80!black] at (0,-1.35) {defect};
    \end{scope}
    \draw[arrow] (-0.9,0)--(0.55,0);
    \node[align=center] at (-0.17,0.52) {equals};
    \begin{scope}[shift={(3.2,0)}]
      \draw[defect] (-1.7,-1.0)--(1.7,-1.0);
      \node[op] (c) at (-0.75,0.82) {};
      \node[op] (d) at (0.78,0.72) {};
      \draw[line] (c)--(-0.75,-1.0);
      \draw[line] (d)--(0.78,-1.0);
      \node[blue!65!black,align=center] at (0,0.12)
        {two\\bulk-to-defect\\OPEs};
      \node[violet!80!black] at (0,-1.35) {defect channel};
    \end{scope}
  \end{tikzpicture}
  }
  \caption{Defect crossing equates the ordinary bulk OPE channel with the
  channel obtained by expanding each bulk operator into defect primaries.}
  \label{fig:defect-crossing-two-channels}
\end{figure}

\section{Boundary CFT}

A boundary CFT is the special case \(p=D-1\).  The defect is a codimension-one
boundary, and the preserved conformal group is \(SO(D,1)\).  A scalar primary
near a flat boundary has the boundary OPE
\[
  \mathcal O(\sigma,y)
  =
  \sum_{\widehat{\mathcal O}}
  b_{\mathcal O\widehat{\mathcal O}}\,
  y^{\widehat\Delta-\Delta}
  \widehat{\mathcal O}(\sigma)+\cdots ,
  \qquad
  y>0,
\]
where \(y\) is the distance to the boundary.  Boundary conditions are
therefore classified by the boundary operator spectrum, the boundary OPE, and
the compatibility of these data with bulk OPE associativity.

The displacement operator is a universal boundary or defect operator.  It is
defined by the Ward identity for broken translations transverse to the defect:
\[
  \partial_\mu T^{\mu i}(x)
  =
  \delta_\Sigma(x_\perp)\,D^i(\sigma),
\]
where \(i\) is a transverse index and \(D^i\) is the displacement operator.
For a \(p\)-dimensional conformal defect, its protected dimension is
\[
  \widehat\Delta_D=p+1.
\]
The displacement two-point coefficient is a central observable of the defect
CFT.

\section{Topological Operators and Higher Symmetries}

Extended operators can also express global symmetries.  A \(q\)-form global
symmetry in \(D\) dimensions is generated by topological operators
\[
  U_g(\Sigma^{D-q-1}),
  \qquad
  g\in G,
\]
supported on closed codimension-\((q+1)\) manifolds.  Topological means that
correlation functions are invariant under smooth deformations of
\(\Sigma^{D-q-1}\) that do not cross charged operators.  A charged
\(q\)-dimensional operator \(V(C^q)\) transforms by linking:
\[
  U_g(\Sigma^{D-q-1})\,V(C^q)
  =
  \chi_R(g)^{\operatorname{Link}(\Sigma,C)}
  V(C^q)
\]
when \(V\) has charge representation \(R\) and the symmetry is abelian or has
been diagonalized in the relevant sector.

\begin{figure}[H]
  \centering
  \resizebox{0.90\textwidth}{!}{%
  \begin{tikzpicture}[x=1cm,y=1cm,every node/.style={font=\scriptsize}]
    \tikzset{
      surface/.style={violet!80!black,line width=1.1pt},
      charged/.style={blue!65!black,line width=1.1pt},
      arrow/.style={-{Latex[length=2mm]},line width=0.85pt}
    }
    \draw[charged] (-3.65,-0.75)
      .. controls (-2.55,0.90) and (-1.30,0.90) .. (-0.35,-0.45);
    \node[blue!65!black] at (-2.15,1.15) {charged operator \(V(C)\)};
    \begin{scope}[shift={(1.05,0)}]
      \draw[surface] (-1.05,0) arc (180:360:1.05 and 0.55);
      \draw[surface,dashed] (1.05,0) arc (0:180:1.05 and 0.55);
      \draw[charged] (0,-1.05)--(0,-0.18);
      \draw[charged,dashed] (0,-0.18)--(0,0.18);
      \draw[charged] (0,0.18)--(0,1.05);
      \node[violet!80!black] at (0,0.95) {topological \(U_g(\Sigma)\)};
      \node[blue!65!black] at (0,-1.35) {linked segment of \(V(C)\)};
    \end{scope}
    \draw[arrow,green!45!black] (2.30,-0.48)--(1.68,-0.12);
    \node[green!45!black,align=center] at (3.15,-0.90)
      {action depends on\\linking data};
  \end{tikzpicture}
  }
  \caption{A higher-form symmetry operator is topological and acts on charged
  extended operators by linking.}
  \label{fig:higher-form-symmetry-linking}
\end{figure}

In a conformal theory, these topological operators are part of the extended
CFT data.  Their fusion rules, correlation functions, anomalies, and actions
on line and surface operators constrain the spectrum of admissible defects and
the possible global forms of gauge theories.

\section{Line Operators and Duality}

Four-dimensional \(\mathcal N=4\) super-Yang--Mills supplies a central
example.  Wilson, 't Hooft, and Wilson-'t Hooft line operators are conformal
line defects when they are supported on a straight line or a circle with the
appropriate scalar coupling.  Their charges form an electric-magnetic lattice
equipped with the Dirac pairing.  The duality group acts on this lattice and
therefore on the set of line defects.

Thus the global form of the gauge group, the allowed line-operator lattice,
and the duality quotient of the conformal manifold are linked pieces of one
extended CFT.  Local correlators, defect correlators, and higher-symmetry
operators must transform together under duality.  This is the bridge from the
local conformal data developed in this volume to the later systematic study of
generalized symmetries, defects, anomalies, and extended operator algebras.
